\begin{question}{0.1}
    (20分) 各向同性线弹性材料的应变能密度函数为 \( W = \frac{1}{2}\sigma_{ij}\varepsilon_{ij} \)。利用广义胡克定律 \( \sigma_{ij} = \lambda \varepsilon_{kk}\delta_{ij} + 2\mu \varepsilon_{ij} \)，证明：
$$
W = \frac{1}{2}\lambda (\varepsilon_{kk})^2 + \mu \varepsilon_{ij}\varepsilon_{ij}
$$

并进一步将其分解为体积变形能和畸变能两部分。
\end{question}

\begin{solution}
    (20分) 各向同性线弹性材料的本构关系：
$$
\sigma_{ij} = \lambda \varepsilon_{kk} \delta_{ij} + 2\mu \varepsilon_{ij}
$$
应变能密度：
$$
W = \frac{1}{2}\sigma_{ij}\varepsilon_{ij}
$$
代入本构关系：
$$
\begin{aligned}
W &= \frac{1}{2} \left( \lambda \varepsilon_{kk} \delta_{ij} + 2\mu \varepsilon_{ij} \right) \varepsilon_{ij} \\
&= \frac{1}{2} \lambda \varepsilon_{kk} \delta_{ij} \varepsilon_{ij} + \frac{1}{2} \cdot 2\mu \varepsilon_{ij} \varepsilon_{ij} \\
&= \frac{1}{2} \lambda \varepsilon_{kk} \varepsilon_{ii} + \mu \varepsilon_{ij} \varepsilon_{ij} \quad (\text{因为 } \delta_{ij}\varepsilon_{ij} = \varepsilon_{ii} = \varepsilon_{kk}) \\
&= \frac{1}{2} \lambda (\varepsilon_{kk})^2 + \mu \varepsilon_{ij} \varepsilon_{ij}
\end{aligned}
$$
证毕。
\end{solution}
